\chapterpageanex{Anexos}
\chapterstar{Anexo 1: Estudio de Viabilidad y Factibilidad}

% Redefinir numeración de secciones para el anexo
\renewcommand{\thesection}{A1.\arabic{section}}
\setcounter{section}{0}

% Cambiar nombre de "Cuadro" a "Tabla"
\renewcommand{\tablename}{Tabla}
% Reiniciar contador de tablas
\setcounter{table}{0}
% Formato: A.1, A.2, A.3...
\renewcommand{\thetable}{A.\arabic{table}}

En el siguiente apartado se presentan las tablas que detallan los recursos esenciales que conforman el estudio de viabilidad y factibilidad del proyecto tecnológico, incluyendo los recursos materiales, tecnológicos, humanos y financieros necesarios para el desarrollo e implementación del sistema de gestión de inventario para la Farmacia Selene, con el propósito de mostrar de manera organizada los elementos requeridos para su elaboración.

% =============================================
% TABLA 1: RECURSOS HUMANOS
% =============================================
\section{Recursos Humanos}
\begin{longtable}{|p{3.8cm}|p{4.5cm}|p{5.2cm}|}

\caption{Recursos Humanos} \label{tab:recursos_humanos} \\

\hline
\rowcolor{gray!20}
\textbf{Puesto} & \textbf{Perfil} & \textbf{Función} \\
\hline
\endfirsthead
\hline
\rowcolor{gray!20}
\textbf{Puesto} & \textbf{Perfil} & \textbf{Función} \\
\hline
\endhead
\hline
\multicolumn{3}{|r|}{\small Continúa en la siguiente página} \\
\hline
\endfoot
\hline
\endlastfoot

Desarrollador Backend &
Conocimientos en programación con Python, desarrollo de APIs REST con FastAPI y manejo de bases de datos relacionales &
Implementar lógica de negocio, estructura de BD y comunicación con frontend e IoT. \\
\hline

Desarrollador Frontend &
Conocimientos en HTML, CSS, JavaScript y React, así como en el uso de bibliotecas de componentes como Material-UI. &
Desarrollar la interfaz web del sistema: vistas, formularios, reportes y alertas. \\
\hline

Desarrollador de Sistemas Embebidos &
Conocimientos en programación en C++, manejo de microcontroladores y protocolos de comunicación como I²C y HTTP. &
Programar el ESP32 para lectura del sensor AHT10 y envío de datos al backend. \\
\hline

Personal de Farmacia &
Conocimiento operativo del inventario, ventas y manejo de medicamentos. &
Registrar ventas, consultar existencias y atender alertas del sistema. \\
\hline

\end{longtable}




% =============================================
% TABLA 2: RECURSOS MATERIALES
% =============================================
\section{Recursos Materiales}

\begin{longtable}{|p{3.5cm}|p{2.5cm}|p{8cm}|}

\caption{Recursos Materiales} \label{tab:recursos_materiales} \\

\hline
\rowcolor{gray!20}
\textbf{Componente} & \textbf{Cantidad} & \textbf{Especificaciones} \\
\hline
\endfirsthead

\hline
\rowcolor{gray!20}
\textbf{Componente} & \textbf{Cantidad} & \textbf{Especificaciones} \\
\hline
\endhead

\hline
\multicolumn{3}{|r|}{\small Continúa en la siguiente página} \\
\hline
\endfoot

\hline
\endlastfoot

Microcontrolador ESP32-C5 & 1 &
Microcontrolador con conectividad Wi-Fi 6 (802.11ax) y Bluetooth Low Energy (BLE), utilizado para la lectura de sensores y comunicación con el servidor. \\
\hline

Sensor AHT10 & 1 &
Sensor digital de temperatura y humedad relativa. Precisión de ±0.3~°C en temperatura y ±2\% en humedad. Se conecta al microcontrolador mediante un protocolo de comunicación digital de corto alcance (I²C). \\
\hline

Lector de códigos de barras & 1 &
Dispositivo de entrada con conexión USB para la captura automática de los códigos de los productos farmacéuticos. \\
\hline

Gabinete & 1 &
Caja de protección utilizada para alojar de forma segura el microcontrolador, el sensor y la fuente de alimentación durante su instalación en la farmacia. \\
\hline

Cables Dupont & 1 juego &
Cables de conexión utilizados para unir físicamente el microcontrolador ESP32 con el sensor AHT10. \\
\hline

Fuente de alimentación USB & 1 &
Adaptador de corriente de 5~V para la alimentación eléctrica del microcontrolador ESP32. \\
\hline

Cable USB & 1 &
Cable utilizado para la programación inicial y alimentación del microcontrolador desde una computadora. \\
\hline

Computadora de desarrollo & 1 &
Equipo de cómputo destinado a ejecutar las herramientas de programación, administración de base de datos y navegador web durante el desarrollo del sistema. \\
\hline

Router / Red Wi-Fi & 1 &
Infraestructura de red local necesaria para la comunicación inalámbrica entre el microcontrolador, el servidor y la interfaz web. \\
\hline

\end{longtable}

\newpage





% =============================================
% TABLA 3: RECURSOS TECNOLÓGICOS
% =============================================
\section{Recursos Tecnológicos}

\begin{longtable}{|p{3.5cm}|p{8cm}|p{2.5cm}|}

\caption{Recursos Tecnológicos} \label{tab:recursos_tecnologicos} \\

\hline
\rowcolor{gray!20}
\textbf{Recurso Tecnológico} & \textbf{Descripción} & \textbf{Tipo} \\
\hline
\endfirsthead

\hline
\rowcolor{gray!20}
\textbf{Recurso Tecnológico} & \textbf{Descripción} & \textbf{Tipo} \\
\hline
\endhead

\hline
\multicolumn{3}{|r|}{\small Continúa en la siguiente página} \\
\hline
\endfoot

\hline
\endlastfoot

Python &
Lenguaje de programación utilizado para el desarrollo del servidor y la lógica de negocio del sistema. &
Software \\
\hline

C++ &
Lenguaje de programación utilizado para programar el microcontrolador ESP32, gestionando la lectura de sensores y el envío de datos al servidor. &
Software \\
\hline

FastAPI &
Herramienta para la construcción del servidor web que gestiona la comunicación entre la interfaz de usuario, los dispositivos físicos y la base de datos. &
Software \\
\hline

PostgreSQL &
Sistema de base de datos relacional utilizado para almacenar la información de usuarios, productos, ventas y registros ambientales del sistema. &
Software \\
\hline

SQLAlchemy &
Herramienta que permite al código Python interactuar con la base de datos sin necesidad de escribir instrucciones SQL complejas directamente. &
Software \\
\hline

React &
Biblioteca de JavaScript utilizada para construir la interfaz visual del sistema accesible desde el navegador web. &
Software \\
\hline

Material-UI (MUI) &
Biblioteca de componentes visuales prediseñados que garantizan consistencia y adaptabilidad en la interfaz web del sistema. &
Software \\
\hline

Arduino IDE / ESP-IDF &
Entorno de programación utilizado para escribir, compilar y cargar el código en el microcontrolador ESP32. &
Software \\
\hline

Vercel &
Plataforma en la nube utilizada para publicar y mantener disponible la interfaz web del sistema. &
Servicio en la nube \\
\hline

Railway &
Plataforma en la nube utilizada para alojar y mantener en operación el servidor y la base de datos del sistema. &
Servicio en la nube \\
\hline

GitHub &
Plataforma de control de versiones utilizada para el almacenamiento, respaldo y gestión del código fuente del proyecto. &
Servicio en la nube \\
\hline

Alembic &
Herramienta de gestión de migraciones de base de datos, utilizada para aplicar y controlar los cambios en la estructura del sistema de almacenamiento de forma ordenada y segura. &
Software \\
\hline

Visual Studio Code &
Entorno de desarrollo utilizado principalmente para la programación del frontend con React y la gestión general de archivos del proyecto. &
Software \\
\hline

PyCharm &
Entorno de desarrollo integrado utilizado para la programación, depuración y estructuración del servidor desarrollado con Python y FastAPI. &
Software \\
\hline

DataGrip &
Herramienta de administración visual de bases de datos utilizada para consultar, validar y gestionar la información almacenada en PostgreSQL. &
Software \\
\hline

JSON (JavaScript Object Notation) &
Formato ligero de intercambio de datos utilizado para transmitir información estructurada entre el frontend, el backend y otros componentes del sistema mediante la API. &
Estándar de datos \\
\hline

JWT (JSON Web Token) &
Mecanismo de autenticación utilizado para verificar la identidad de los usuarios y controlar el acceso seguro a los endpoints de la API mediante tokens firmados digitalmente. &
Protocolo de seguridad \\
\hline

\end{longtable}
\newpage







% =============================================
% TABLA 4: RECURSOS FINANCIEROS 
% =============================================
\section{Recursos Financieros}


% ===============================
% COSTOS HARDWARE
% ===============================
\pgfmathsetmacro{\ESPPrecio}{220}
\pgfmathsetmacro{\SensorPrecio}{150}
\pgfmathsetmacro{\ScannerPrecio}{500}
\pgfmathsetmacro{\CablesPrecio}{50}
\pgfmathsetmacro{\FuentePrecio}{120}
\pgfmathsetmacro{\USBPrecio}{50}
\pgfmathsetmacro{\GabinetePrecio}{120}

\pgfmathsetmacro{\HardwareTotal}{%
  \ESPPrecio + \SensorPrecio + \ScannerPrecio +
  \CablesPrecio + \FuentePrecio + \USBPrecio + \GabinetePrecio
}

% ===============================
% COSTOS CLOUD
% ===============================
\pgfmathsetmacro{\VercelMensual}{360}
\pgfmathsetmacro{\RailwayMensual}{360}
\pgfmathsetmacro{\VercelAnual}{\VercelMensual * 12}
\pgfmathsetmacro{\RailwayAnual}{\RailwayMensual * 12}
\pgfmathsetmacro{\MensualTotal}{\VercelMensual + \RailwayMensual}
\pgfmathsetmacro{\AnualTotal}{\MensualTotal * 12}

\begin{center}
\begin{longtable}{
  | p{3cm}      % Componente
  | p{2cm}      % Precio Unitario
  | p{1.5cm}    % Cantidad
  | p{2cm}      % Total Inicial
  | p{2cm}      % Costo Mensual
  | p{2cm} |    % Costo Anual
}

\caption{Recursos Financieros} \label{tab:rec_financieros} \\
\hline
\rowcolor{gray!20}
\textbf{Componente} &
\makecell{\textbf{Precio}\\\textbf{Unitario}\\\textbf{(MXN)}} &
\makecell{\textbf{Cant.}} &
\makecell{\textbf{Total}\\\textbf{Inicial}\\\textbf{(MXN)}} &
\makecell{\textbf{Costo}\\\textbf{Mensual}\\\textbf{(MXN)}} &
\makecell{\textbf{Costo}\\\textbf{Anual}\\\textbf{(MXN)}} \\
\hline
\endfirsthead

\hline
\rowcolor{gray!20}
\textbf{Componente} &
\makecell{\textbf{Precio}\\\textbf{Unitario}\\\textbf{(MXN)}} &
\makecell{\textbf{Cant.}} &
\makecell{\textbf{Total}\\\textbf{Inicial}\\\textbf{(MXN)}} &
\makecell{\textbf{Costo}\\\textbf{Mensual}\\\textbf{(MXN)}} &
\makecell{\textbf{Costo}\\\textbf{Anual}\\\textbf{(MXN)}} \\
\hline
\endhead

\hline
\multicolumn{6}{|r|}{\small Continúa en la siguiente página} \\
\hline
\endfoot

\hline
\endlastfoot

% ===============================
% HARDWARE
% ===============================
ESP32-C5 Microcontrolador & \num{\ESPPrecio} & 1 & \num{\ESPPrecio} & 0 & 0 \\
\hline
Sensor AHT10 & \num{\SensorPrecio} & 1 & \num{\SensorPrecio} & 0 & 0 \\
\hline
Lector de códigos de barras & \num{\ScannerPrecio} & 1 & \num{\ScannerPrecio} & 0 & 0 \\
\hline
Cables Dupont & \num{\CablesPrecio} & 1 & \num{\CablesPrecio} & 0 & 0 \\
\hline
Fuente de alimentación USB & \num{\FuentePrecio} & 1 & \num{\FuentePrecio} & 0 & 0 \\
\hline
Cable USB & \num{\USBPrecio} & 1 & \num{\USBPrecio} & 0 & 0 \\
\hline
Gabinete & \num{\GabinetePrecio} & 1 & \num{\GabinetePrecio} & 0 & 0 \\
\hline

% ===============================
% SERVICIOS CLOUD
% ===============================
Vercel (Hosting Frontend, Pro) & \num{360} & 1 & 0 & \num{\VercelMensual} & \num{\VercelAnual} \\
\hline
Railway (Hosting Backend + Base de Datos) & \num{360} & 1 & 0 & \num{\RailwayMensual} & \num{\RailwayAnual} \\
\hline

% ===============================
% TOTALES
% ===============================
\multicolumn{3}{|r|}{\textbf{Total inversión inicial}} & \num{\HardwareTotal} & & \\
\hline
\multicolumn{4}{|r|}{\textbf{Costo mensual del sistema}} & \num{\MensualTotal} & \\
\hline
\multicolumn{5}{|r|}{\textbf{Costo anual del sistema}} & \num{\AnualTotal} \\
\hline

\end{longtable}
\end{center}
\newpage

\section{Impacto Social}

La implementación del sistema de gestión de inventario para la Farmacia Selene responde a la necesidad de sustituir un proceso manual propenso a errores, pérdidas económicas y falta de información actualizada. Su digitalización optimiza la operación interna del negocio y genera beneficios directos para la comunidad.

En primer lugar, mejora la disponibilidad y confiabilidad de los medicamentos mediante un control de inventario centralizado y actualizado, reduciendo el desabasto y la acumulación de productos caducados. Esto garantiza un acceso más oportuno y seguro a los productos farmacéuticos.

En segundo lugar, contribuye a la conservación de la calidad de los medicamentos al integrar sensores de temperatura y humedad que permiten supervisar las condiciones de almacenamiento. Este monitoreo continuo mantiene buena actualización de los productos en buen estado y favorece el cumplimiento de buenas prácticas farmacéuticas.

Asimismo, el sistema fortalece al personal al disminuir tareas manuales repetitivas y permitir una mayor concentración en la atención al cliente, promoviendo además la adopción de herramientas tecnológicas en el entorno laboral.

Finalmente, el proyecto demuestra que soluciones basadas en aplicaciones web, bases de datos en la nube y monitoreo ambiental son viables en pequeños establecimientos, sentando un precedente para la modernización de negocios similares en el ámbito local.

